\fancyhead[L]{\small LETTER TO BRANDES.}
\SourceLeaf{515}\pdfbookmark[0]{Letter from Gauss to Brandes}{brief_brandes}
\expandafter\def\csname HistoricalNoteText515\endcsname{[On Achromatic Doublet Objectives, with Particular Regard to the Complete Elimination of Colour Dispersion.]}
\begin{center}
\hrule\vspace{2mm}
\textsc{Gehler's} Physikalisches Wörterbuch. 1831. Article:\par
Lenses, Calculations on Achromatic and Aplanatic Objectives Made of Two Glass Lenses.
\vspace{2mm}\hrule\vspace{5mm}
{\large\itshape Letter from Gauss to Brandes.\par}
\end{center}
Your letter prompted me to spend a spare hour on the point briefly mentioned at the end of that paper\HistoricalNote{515}. The real meaning of the remark there does not seem to have been understood quite correctly by everyone, and my own statement also needs a slight modification. I now find, through a \textit{more searching} investigation, that the blur which depends on the fourth power of the distance of the incident rays from the axis in the expression for longitudinal aberration has the smallest possible overall effect if the objective is constructed so that the rays that strike infinitely close to the axis and those that would strike it at a distance $=R.\BareRoot{\tfrac65}$ (where $R=$ the radius of the objective) meet at \textit{one} point $A$. The eyepiece is then set to give a clear view of the point on the axis where all the rays that struck at distances $=(\tfrac35-\tfrac{\BareRoot{6}}{10})R$ and $=(\tfrac35+\tfrac{\BareRoot{6}}{10})R$ from the axis meet. Imagine a plane through this point perpendicular to the axis. The image is then the less clear, the larger the circle around $A$ filled by the rays that strike the objective from a point of the object, provided that the intensity of the rays at each place in the circle is also taken into account. There is some arbitrariness here. I think it best to proceed on the same principles that underlie the method of least squares. If $\mathrm{d}s$ is an element of this circle, $\rho$ the distance of the element from $A$, and $i$ the intensity of the rays there, I assume that $\int i\rho\rho\,\mathrm{d}s$ can be taken as the measure of the total blur\SourceLeaf{516} and minimize it. I find the following results. 1. Suppose the objective is constructed so that the term in the longitudinal aberration that depends on the square of the distance from the axis becomes $=0$, and the eyepiece is set so. that $A$ lies where the rays infinitely close to the axis intersect it. Let the value of this integral be $=E$. 2. If, however, \textit{with the same arrangement}, the eyepiece is set to make the integral as small as possible for that arrangement (in which case $A$ will be the point where the rays striking at a distance $=R\BareRoot[4]{\tfrac12}$ meet), the integral is $=\tfrac14E$. 3. With the arrangement described above and the best setting of the eyepiece, by contrast, the integral is $=\tfrac{1}{100}E$, the absolute minimum. The result above, that a merely \textit{hypothetical} image (from rays farther from the axis than the radius of the objective) should be made to coincide with the point where the rays infinitely close to the axis meet, is at first very surprising and seems paradoxical. On closer consideration, however, the actual reason is easy to see. That first \textit{so-called} principal image (from rays very close to the axis) is, so to speak, the least important here because of its low intensity. It is much more important to keep the rays from the rings nearer the edge of the objective closer to \textit{one another}, which that arrangement does best. I regret that the limits of a letter do not allow me to go into more detail now. There is no disputing a precise calculation, and in vague reasoning one easily overlooks an essential point. These hints will suffice for a reader familiar with the subject.

In general I find that, \textit{with the best setting of the eyepiece}, the integral is always $=\tfrac14E(1-\tfrac85\mu\mu+\tfrac23\mu^4)$ if the objective is constructed so that the rays at a distance $\mu R$ from the axis meet at one point with the image called the principal image above. It is a minimum for $\mu=\BareRoot{\tfrac65}$ and is then $=\tfrac{1}{100}E$. For $\mu=1$ it would only be $=\tfrac{1}{60}E$, and for $\mu=$ infinitely small, $=\tfrac14E$. It follows that \textsc{Bohnenberger} was mistaken and that I too was mistaken at the time, in that I had still not departed far enough from \textsc{Bohnenberger}. I had then considered only the \textit{overall size} of the blurred image, without taking into account the unequal intensity of its individual parts.

